\documentclass[11pt,a4paper]{article}
\usepackage{vespa_style}

\begin{document}

\vespatitlepage
  {Rapport technique}
  {Guide de reproduction — matériel, système, firmware}
  {Document de clôture — juillet 2026\\
   Tout ce qui est nécessaire pour reconstruire le prototype mk.2 depuis zéro.}

\tableofcontents
\newpage

% =====================================================================
\section{Objet de ce document}

Ce document est un \textbf{guide de reproduction}. Il décrit, dans l'ordre où il faut les
exécuter, les opérations permettant de reconstituer le prototype \vespa{} mk.2 :
approvisionnement, câblage, installation système, compilation des deux firmwares,
calibration, mise en route.

Il ne justifie pas les choix de conception — cela relève du rapport de soutenance
(\fichier{docs/Rapport\_PFE\_2026\_VESPA.pdf}). Il ne raconte pas non plus les impasses —
cela relève de \fichier{docs/Difficultes\_techniques.pdf}.

\begin{avertissement}[Lisez d'abord « Difficultés techniques »]
Plusieurs étapes de ce guide comportent des pièges qui \textbf{ne se voient pas} et dont le
symptôme n'a aucun rapport apparent avec la cause (une broche qui coupe la console série,
un outil NVIDIA qui détruit la configuration des caméras). Elles sont signalées ici, mais
elles ne sont \emph{expliquées} que dans le document des difficultés. Le temps passé à le
lire sera regagné plusieurs fois.
\end{avertissement}

\paragraph{État de ce qui est reproductible.} Ce guide mène à un système qui \textbf{suit
une cible ArUco en 3D et oriente la tourelle}. Il ne mène pas à un système qui reconnaît un
frelon (modèle indisponible), ni qui tire (watchdog non implémenté, laser jamais mis sous
tension). Voir la synthèse en fin de document.

% =====================================================================
\section{Architecture matérielle}

\begin{center}
\begin{tabularx}{\textwidth}{@{}llX@{}}
\toprule
\textbf{N\oe ud} & \textbf{Calculateur} & \textbf{Rôle} \\
\midrule
\texttt{vision\_node} & Jetson Orin Nano Super 8~Go &
Acquisition stéréo synchronisée, détection, triangulation, prédiction, émission de la
consigne \\
\addlinespace
\texttt{motion\_node} & NUCLEO-G431RB &
Asservissement FOC en boucle fermée des axes \textsc{pan} et \textsc{tilt} \\
\addlinespace
\texttt{coordination\_node} & ESP32 &
Collimation, tir, watchdog — \textcolor{vespared}{\textbf{non implémenté}} \\
\bottomrule
\end{tabularx}
\end{center}

\noindent
Les n\oe uds communiquent sur un \textbf{bus CAN 2.0B à 500~kbit/s} (trames standard
11~bits). Le dictionnaire de données complet est en annexe
(\fichier{docs/annexes/protocole\_can.md}) ; sa définition normative est
\fichier{software/common/can\_protocol.h}.

% =====================================================================
\section{Nomenclature (BOM)}

Nomenclature du prototype mk.2, telle qu'effectivement approvisionnée.
Total $\approx$ \textbf{830~\euro{}} TTC.

\begin{center}
\begin{tabularx}{\textwidth}{@{}Xlr@{}}
\toprule
\textbf{Élément} & \textbf{Qté} & \textbf{Prix TTC} \\
\midrule
\multicolumn{3}{@{}l}{\textit{\textcolor{vespagray}{Calcul et électronique}}} \\
Jetson Orin Nano Super Developer Kit & 1 & 311,18~\euro{} \\
B-G431B-ESC1 \textit{\footnotesize (voir note ci-dessous)} & 2 & 16,31~\euro{} \\
ESP32-C3 SuperMini \textit{\footnotesize (coordination — non utilisé)} & 1 & 4,29~\euro{} \\
Alimentation AC-DC 24~V / 150~W & 1 & 5,29~\euro{} \\
Transceivers CAN (lot de 10, type WCMCU-230 / SN65HVD230) & 1 & 5,19~\euro{} \\
TMC2209 (lot de 4) \textit{\footnotesize (collimation)} & 1 & 9,09~\euro{} \\
\addlinespace
\multicolumn{3}{@{}l}{\textit{\textcolor{vespagray}{Vision}}} \\
Module caméra Arducam OV9281 1~MP monochrome \emph{global shutter} & 2 & 142,59~\euro{} \\
Objectif M12 5~MP 4,6~mm sans distorsion & 2 & 9,29~\euro{} \\
\addlinespace
\multicolumn{3}{@{}l}{\textit{\textcolor{vespagray}{Mécanique et actionneurs}}} \\
iFlight iPower GM5208-12 + encodeur AS5048A (paire) & 1 & 134,99~\euro{} \\
Actionneur linéaire course 35~mm & 1 & 20,42~\euro{} \\
Actionneur linéaire course 10~mm & 2 & 1,99~\euro{} \\
Capteur photoélectrique en U \textit{\footnotesize (prise d'origine)} & 1 & 2,54~\euro{} \\
\addlinespace
\multicolumn{3}{@{}l}{\textit{\textcolor{vespagray}{Optique de tir}}} \\
\textbf{Diode laser bleue 5,5~W} \textcolor{vespared}{\textbf{(Classe 4)}} & 1 & 83,69~\euro{} \\
Lentille asphérique K9 D63 F100 \textit{\footnotesize (L2, focale 100~mm)} & 1 & 23,59~\euro{} \\
\addlinespace
\multicolumn{3}{@{}l}{\textit{\textcolor{vespagray}{Sécurité — non optionnel}}} \\
\textbf{Lunettes de protection OD7+ CE 190--550~nm} & 1 & 30,29~\euro{} \\
\bottomrule
\end{tabularx}
\end{center}

\begin{note}[Deux divergences entre la nomenclature et le matériel réellement utilisé]
\begin{enumerate}[nosep]
  \item \textbf{Cartes de commande moteur.} La nomenclature liste des \textbf{B-G431B-ESC1}.
        Le firmware livré cible en réalité une \textbf{NUCLEO-G431RB} portant deux shields
        \textbf{X-NUCLEO-IHM16M1} (STSPIN830) — c'est ce que déclarent
        \fichier{platformio.ini} (\texttt{board = nucleo\_g431rb}) et
        \fichier{pinout.h}. \textbf{Suivre le code, pas la nomenclature.}
  \item \textbf{Le laser.} La nomenclature (et donc le matériel acheté) porte sur une diode
        \textbf{5,5~W bleue}. Le dimensionnement théorique du rapport de soutenance a été
        conduit pour $P = 5$~W à 460~nm. Les deux valeurs sont cohérentes en ordre de
        grandeur ; c'est la diode 5,5~W qui est physiquement présente.
        \textbf{Classe 4 dans les deux cas.}
\end{enumerate}
\end{note}

\paragraph{Moteurs.} GM5208-12, PMSM \emph{pancake} : \textbf{7 paires de pôles}
(14 aimants), $R_s = 7{,}48~\Omega$, $L_s = 4{,}49$~mH, tension nominale 24~V, courant
nominal 1,5~A. Encodeur magnétique absolu AS5048A, \textbf{14 bits} $\rightarrow$ résolution
$360/2^{14} \approx 0{,}022^\circ$, borne théorique de précision de la tourelle.

% =====================================================================
\section{Câblage}

Le câblage complet, broche à broche, est en annexe :
\textbf{\fichier{docs/annexes/cablage.md}}. Il n'est pas reproduit ici, mais
\textbf{deux points sont critiques} et doivent être connus avant de souder :

\begin{avertissement}[Les deux pièges du câblage]
\textbf{1. Le potentiomètre \textsc{pan} n'est PAS sur \broche{PA2}, mais sur
\broche{PB14}.}\\
\broche{PA2} est câblée sur \texttt{USART2\_TX} (console série du ST-LINK) sur la
NUCLEO-G431RB. Y brancher le potentiomètre \textbf{coupe la console}. Le fichier historique
\fichier{STM32\_WiringManifest.txt} indique encore \broche{PA2} : il est faux.

\medskip
\textbf{2. \broche{PA4} est réservée au SPI (NSS matériel), pas au retour de courant.}\\
Elle doit rester haute, sous peine de faire basculer le SPI1 — qui porte les deux encodeurs
— en mode esclave. Le retour de courant de la phase W du moteur 1 est sacrifié en
conséquence.
\end{avertissement}

\paragraph{Bus CAN.} Transceiver \textbf{3,3~V uniquement} (SN65HVD230 / WCMCU-230).
\textbf{Ne jamais l'alimenter en 5~V} : cela détruit le composant et peut endommager les
broches du Jetson.
\begin{itemize}[nosep]
  \item Côté Jetson : \texttt{CTX} $\rightarrow$ broche \textbf{31}, \texttt{CRX}
        $\rightarrow$ broche \textbf{29}, 3V3 $\rightarrow$ broche 1, GND $\rightarrow$
        broche 6.
  \item Côté STM32 : \texttt{CAN\_TX} = \broche{PB9}, \texttt{CAN\_RX} = \broche{PB8}.
  \item \textbf{Masse commune obligatoire} entre les n\oe uds (3\ieme{} fil) : le CAN est
        différentiel, mais le transceiver a une limite de tension de mode commun.
  \item Terminaison \textbf{120~$\Omega$} aux \emph{deux extrémités} du bus, et
        \emph{seulement} aux extrémités. Les modules WCMCU-230 embarquent une résistance
        (sérigraphiée \texttt{R2}, marquage « 121 ») : la dessouder sur tout n\oe ud qui
        n'est pas en bout de bus.
\end{itemize}

% =====================================================================
\newpage
\section{Installation du Jetson Orin Nano}

\begin{avertissement}[L'étape la plus délicate du projet]
Le Jetson ne fonctionne \textbf{pas} avec le noyau JetPack standard : les caméras OV9281
exigent le \textbf{noyau Arducam}. Et une fois cette configuration en place, l'outil
\fichier{jetson-io.py} \textbf{ne doit plus jamais être exécuté} — il régénère globalement
les \emph{overlays} d'arbre de périphériques et \textbf{détruit la synchronisation stéréo}.
\end{avertissement}

\subsection{Système de base}

\begin{enumerate}
  \item Flasher le \textbf{JetPack 6.2} (L4T r36.4) sur le Jetson Orin Nano Super, via
        NVIDIA SDK Manager. Racine sur NVMe recommandée (c'est la configuration d'origine).
  \item Installer le \textbf{noyau et les pilotes Arducam} pour l'OV9281, en suivant la
        procédure du constructeur. Cela installe un noyau dédié dans
        \fichier{/boot/arducam/}.
\end{enumerate}

\subsection{Overlays d'arbre de périphériques}

La configuration matérielle est portée par trois \emph{overlays} chargés au démarrage. Les
fichiers de référence sont conservés dans \fichier{jetson/} :

\begin{center}
\begin{tabularx}{\textwidth}{@{}lX@{}}
\toprule
\textbf{Overlay} & \textbf{Rôle} \\
\midrule
\fichier{...-arducam-dual.dtbo} &
Double caméra OV9281. C'est lui qui résout le conflit apparent d'adresse
I\textsuperscript{2}C (un bus par port CSI). \\
\addlinespace
\fichier{jetson-io-hdr40-user-custom.dtbo} &
\textbf{PWM de synchronisation} (broches 32 et 33 de l'en-tête 40 points).
\textcolor{vespared}{\textbf{Ne pas modifier ni régénérer.}} \\
\addlinespace
\fichier{can0-pinmux.dtbo} &
Tentative de routage du CAN0. \textcolor{vespared}{\textbf{Ne fonctionne pas}} — voir
« Difficultés techniques ». \\
\bottomrule
\end{tabularx}
\end{center}

\noindent
Le fichier \fichier{/boot/extlinux/extlinux.conf} de référence est fourni
(\fichier{jetson/extlinux.conf.reference}). L'entrée active y charge le noyau Arducam et
la liste d'\emph{overlays} :

\begin{lstlisting}[language=bash]
LABEL JetsonIO
    MENU LABEL Custom Header Config: <CSI Camera ARDUCAM Dual>
    LINUX /boot/arducam/Image
    FDT   /boot/arducam/dts/dtb/tegra234-p3768-0000+p3767-0005-nv-super.dtb
    INITRD /boot/initrd
    APPEND ${cbootargs} root=PARTUUID=... rw rootwait rootfstype=ext4 ...
    OVERLAYS /boot/arducam/dts/tegra234-p3767-camera-p3768-arducam-dual.dtbo,
             /boot/jetson-io-hdr40-user-custom.dtbo,
             /boot/can0-pinmux.dtbo
\end{lstlisting}

\subsection{Vérifications}

\begin{lstlisting}[language=bash]
# Les deux cameras doivent apparaitre
ls /dev/video*                       # -> /dev/video0  /dev/video1
v4l2-ctl --list-devices

# Le PWM de synchronisation doit etre exportable
ls /sys/class/pwm/pwmchip3/

# Le CAN physique : NE FONCTIONNE PAS a ce jour (verrou ouvert)
sudo bash -c "cat /sys/kernel/debug/pinctrl/*/pinmux-pins | grep -i paa"
#   -> aucune sortie = les broches PAA ne sont pas reclamees
\end{lstlisting}

\subsection{Bus CAN virtuel (indispensable pour travailler)}

Le bus physique étant inopérant, le développement et la validation se font sur
\texttt{vcan0}. Le fichier \fichier{jetson/systemd/vcan.conf} charge le module au
démarrage :

\begin{lstlisting}[language=bash]
sudo modprobe vcan
sudo ip link add dev vcan0 type vcan
sudo ip link set up vcan0
candump vcan0                        # observer les trames
\end{lstlisting}

\noindent
Le service \fichier{jetson/systemd/can0-setup.service} configure le bus \emph{physique}
(500~kbit/s) — il est écrit pour ne se déclencher que si le noyau crée réellement le
périphérique \texttt{can0}, ce qui n'est jamais arrivé.

% =====================================================================
\newpage
\section{Compilation}

\subsection{N\oe ud de vision (Jetson)}

\paragraph{Dépendances.}
\begin{lstlisting}[language=bash]
sudo apt install libopencv-dev libi2c-dev libzmq3-dev cppzmq-dev \
                 v4l-utils can-utils cmake build-essential
\end{lstlisting}

\begin{note}[Dépendances exigées mais non utilisées]
Le \fichier{CMakeLists.txt} déclare \texttt{find\_package(CUDA REQUIRED)} et
\texttt{find\_package(VPI 3.0 REQUIRED)}, ainsi qu'un test sur \fichier{i2c/smbus.h}.
\textbf{Le chemin de code actuel (ArUco + OpenCV) n'en a pas besoin} : ces dépendances sont
un vestige de la voie TensorRT abandonnée. Elles restent \emph{obligatoires} pour que
CMake aboutisse. Si la reprise ne vise pas l'inférence, ces trois exigences peuvent être
retirées sans dommage.
\end{note}

\paragraph{Construction.}
\begin{lstlisting}[language=bash]
cd software/vision_node
mkdir -p build && cd build
cmake .. && make -j6
\end{lstlisting}

\paragraph{Cibles produites.}
\begin{center}
\begin{tabularx}{\textwidth}{@{}lX@{}}
\toprule
\textbf{Binaire} & \textbf{Usage} \\
\midrule
\texttt{vision\_node} & \textbf{Cible de production.} Sans interface, sans ZeroMQ. Émet sur le CAN. \\
\texttt{stereo\_tracker} & Idem + flux ZeroMQ pour la visualisation (\emph{live viewport}). \\
\texttt{calibration\_bridge} & Pont ZeroMQ pour les outils Python de calibration. \\
\texttt{test\_camera\_hal} & Test unitaire d'acquisition. \\
\texttt{test\_hardware\_sync} & \textbf{Test de la synchronisation stéréo} — à lancer en premier. \\
\texttt{test\_stereo\_pipeline} & Test de la triangulation hors ligne. \\
\bottomrule
\end{tabularx}
\end{center}

\paragraph{Exécution.} Les droits \emph{root} sont nécessaires (accès à
\fichier{/sys/class/pwm} et \fichier{/dev/video*}) :
\begin{lstlisting}[language=bash]
sudo ./vision_node
\end{lstlisting}

\subsection{N\oe ud de mouvement (STM32)}

Projet PlatformIO. La bibliothèque commune est liée par lien symbolique, ce qui garantit
qu'il n'existe \textbf{qu'une seule définition} du protocole CAN :

\begin{lstlisting}[language=bash]
cd software/motion_node
pio run                  # compiler
pio run -t upload        # televerser (sonde ST-LINK integree a la Nucleo)
pio device monitor -b 115200
\end{lstlisting}

\noindent
Configuration (\fichier{platformio.ini}) : \texttt{board = nucleo\_g431rb}, framework
Arduino, SimpleFOC 2.3.2, avec activation explicite des modules HAL
\texttt{OPAMP} et \texttt{FDCAN}.

\begin{avertissement}[Ne jamais instancier \texttt{MotionCan} sans arguments]
Le constructeur déclare des valeurs par défaut \texttt{(PA11, PA12)} qui sont
\textbf{fausses} pour ce câblage — \broche{PA11} est le \texttt{EN/FAULT} du moteur~1.
Toujours passer les broches explicitement :
\texttt{MotionCan canBus(CAN\_RX, CAN\_TX);} soit \broche{PB8}/\broche{PB9}.
\end{avertissement}

% =====================================================================
\newpage
\section{Calibration}

Trois outils web (Flask) tournent sur le Jetson et s'utilisent depuis un navigateur
distant — la machine étant sans écran.

\begin{avertissement}[Un seul programme à la fois sur le matériel]
Le n\oe ud de vision et les outils Python se disputent le bus I\textsuperscript{2}C. Un
verrou d'exclusion (\fichier{/tmp/vespa\_hardware.lock}) fait \textbf{échouer volontairement}
le second démarrage. \textbf{Ce n'est pas un bug} : arrêter l'un avant de lancer l'autre.
\end{avertissement}

\subsection{Ordre des opérations}

\begin{enumerate}
  \item \textbf{Exposition et gain} — \fichier{software/tools/calib\_exposition/}
        \begin{lstlisting}[language=bash]
cd software/tools/calib_exposition && ./run_exposure.sh   # http://<jetson>:5000
        \end{lstlisting}
        Régler l'exposition et le gain jusqu'à obtenir un marqueur net et contrasté.
        Les profils sont enregistrés dans
        \fichier{software/vision\_node/config/profiles/} (P1, P2~«~sombre~», P3~«~lumineux~»),
        le profil actif étant désigné par \fichier{config/.active\_profile}.

        \begin{note}[Limite connue]
        Le capteur n'applique pas le même calcul de temps-ligne en mode maître et en mode
        déclenchement externe. \textbf{Les valeurs réglées par cet outil ne correspondent pas
        exactement à celles du \emph{pipeline} C++.} Prévoir un ajustement final en
        conditions réelles.
        \end{note}

  \item \textbf{Calibration stéréo} — \fichier{software/tools/calib\_stereo/}
        \begin{lstlisting}[language=bash]
cd software/tools/calib_stereo && ./run_calibration.sh    # http://<jetson>:5000
        \end{lstlisting}
        Méthode de Zhang \cite{zhang2000calibration} sur mire d'échiquier.
        \textbf{Couvrir tout le champ} — coins, bords, près et loin, et pas seulement le
        centre : sans couverture des bords, la distorsion ne peut pas être résolue, la
        focale « hallucine » et l'erreur RMS explose. Résultat écrit dans
        \fichier{config/stereovision\_settings.json}.

  \item \textbf{Vérification} — \fichier{software/tools/viewport\_live/}
        \begin{lstlisting}[language=bash]
cd software/tools/viewport_live && ./run_liveviewport.sh
        \end{lstlisting}
        Affiche le flux avec superposition : position courante du marqueur et
        \textbf{visée prédictive} (réticule cyan) issue du filtre de Kalman.
\end{enumerate}

% =====================================================================
\section{Mise en route}

\begin{enumerate}
  \item \textbf{Vérifier la synchronisation} avant toute chose :
        \begin{lstlisting}[language=bash]
sudo ./build/test_hardware_sync
        \end{lstlisting}
        La dérive inter-caméras doit être nulle. Si elle ne l'est pas, ne pas aller plus
        loin : la triangulation serait fausse.

  \item \textbf{Monter le bus} (virtuel, en l'état actuel) : \texttt{vcan0}, et lancer
        \texttt{candump} sur une seconde console.

  \item \textbf{Démarrer le n\oe ud de mouvement.} À la mise sous tension, la séquence FOC
        s'initialise. Les décalages d'encodeur sont \textbf{codés en dur} dans
        \fichier{main.cpp} :
        \texttt{motorPan.sensor\_offset = 4{,}7500}, \texttt{motorTilt.sensor\_offset =
        4{,}8179} — ils repoussent le passage $0/2\pi$ hors de la plage utile.
        \textbf{Ils sont propres à un montage mécanique donné} et devront être réétalonnés
        sur un nouveau montage (\fichier{Examples/OneTimeInitialisationAngle.cpp}).

  \item \textbf{Démarrer le n\oe ud de vision} : \texttt{sudo ./vision\_node}.

  \item \textbf{Présenter le marqueur ArUco} dans la \emph{killbox}. La tourelle doit
        suivre.
\end{enumerate}

\paragraph{Butées logicielles} (\fichier{main.cpp}, en radians) — elles protègent la
mécanique et doivent être revues sur un nouveau montage :
\begin{center}
\begin{tabular}{@{}lrr@{}}
\toprule
\textbf{Axe} & \textbf{Min} & \textbf{Max} \\
\midrule
\textsc{pan}  & $-0{,}52$ rad ($-30^\circ$) & $+0{,}52$ rad ($+30^\circ$) \\
\textsc{tilt} & $-0{,}26$ rad ($-15^\circ$) & $+0{,}78$ rad ($+45^\circ$) \\
\bottomrule
\end{tabular}
\end{center}

\paragraph{Arrêt d'urgence.} Bouton bleu (\broche{PC13}). L'arrêt est \textbf{latché} :
coupure des deux étages de puissance, diffusion d'un \emph{heartbeat} \texttt{FAULT}, puis
boucle infinie. \textbf{Seul un reset matériel (bouton noir) permet de repartir.} C'est
délibéré : sur un système portant un laser de Classe 4, aucun redémarrage automatique n'est
acceptable.

% =====================================================================
\newpage
\section{Zone opérationnelle et dimensionnement}

\paragraph{La \emph{killbox}.} Volume dans lequel une cible peut être neutralisée :
parallélépipède de \textbf{500~mm $\times$ 310~mm $\times$ 300~mm}
(longueur $\times$ largeur $\times$ hauteur). La tourelle est placée \textbf{en surplomb},
de sorte que le tir soit dirigé vers le sol.

\paragraph{Optique.} Le compromis porte sur le GSD (\emph{Ground Sampling Distance}) : il ne
doit pas dépasser $\approx 0{,}5$~mm/px pour rester exploitable par un modèle de détection.
Pour un capteur $1\,280 \times 800$, le script
\fichier{simulation/vision/Camera\_PosResAttitude\_v6.m} donne :
\[
  f = 6~\text{mm}, \qquad d_n = 990~\text{mm}
\]

\paragraph{Cadence.} La boucle complète doit tenir dans \textbf{16,6~ms} (60~Hz). À 60~fps,
le système suit une cible se déplaçant jusqu'à $0{,}6~\text{m}\cdot\text{s}^{-1}$ — ce qui
suppose un \textbf{frelon en vol de guet quasi stationnaire}, comportement effectivement
documenté devant les ruches \cite{perrard2009colonie}. Un frelon en vol rapide
($\approx 6~\text{m}\cdot\text{s}^{-1}$) exigerait 600~fps : \textbf{hors de portée}, et
c'est une limite assumée du système.

\paragraph{Obturateur global : pourquoi il est non négociable.} Le flou de mouvement vaut
$b = v \cdot t_{\text{exp}} / R$. À $0{,}6~\text{m}\cdot\text{s}^{-1}$, avec 2~ms
d'exposition et $R = 0{,}5$~mm/px, on obtient $b = 2$~px — sur une cible qui n'en mesure que
20 à 35. Un obturateur à balayage (\emph{rolling shutter}) est donc éliminé. Or aucun
capteur \textbf{trichrome} à obturateur global n'existe à prix abordable : \textbf{d'où le
capteur monochrome, et d'où l'incompatibilité avec VespAI}
(voir « Difficultés techniques »).

% =====================================================================
\section{Structure du dépôt}

\begin{lstlisting}
project_vespa/
|- docs/                  Rapports (PDF + sources LaTeX), annexes, references
|  |- src/                Sources LaTeX + vespa.bib + figures
|  |- annexes/            Cablage, protocole CAN, BOM, setup Jetson, securite
|  \- references/         Articles, datasheets, etudes preparatoires
|- hardware/              Mecanique (impressions 3D), moteurs
|- jetson/                Configuration systeme (overlays, services CAN, extlinux)
|- simulation/            Scripts MATLAB de dimensionnement (vision + optique)
|- software/
|  |- common/             Protocole CAN partage (C, portable STM32/ESP32/Linux)
|  |- vision_node/        Jetson - C++17 / OpenCV / V4L2 / SocketCAN (CMake)
|  |- motion_node/        STM32G431RB - SimpleFOC / FDCAN (PlatformIO)
|  |- coordination_node/  ESP32 - NON IMPLEMENTE (specification seule)
|  |- tools/              Calibration stereo, exposition, viewport
|  \- legacy_mk1_bench/   Banc de caracterisation mk.1 (Arduino Uno)
\- tools/                 Scripts utilitaires du depot
\end{lstlisting}

\paragraph{Le banc mk.1.} \fichier{software/legacy\_mk1\_bench/} contient le code Arduino
ayant servi à caractériser le prototype de première génération (comparaison des algorithmes
« classique » et « accumulateur », mesure du \emph{jitter}, test d'endurance). Ce sont les
\textbf{annexes A à D du rapport de soutenance}. Il est conservé parce qu'il porte les
mesures qui justifient l'abandon des moteurs pas-à-pas.

% =====================================================================
\section{Synthèse de l'état livré}

\begin{center}
\begin{tabularx}{\textwidth}{@{}lXl@{}}
\toprule
\textbf{Sous-système} & \textbf{État} & \textbf{Validé sur} \\
\midrule
Acquisition stéréo & Synchronisation matérielle 60~Hz, dérive 0,000~ms &
\textcolor{vespagreen}{matériel réel} \\
Triangulation & Stéréo éparse sur marqueur ArUco &
\textcolor{vespagreen}{matériel réel} \\
Prédiction & Filtre de Kalman, compense la latence de la chaîne &
\textcolor{vespagreen}{matériel réel} \\
Asservissement & FOC boucle fermée, 2 axes, butées, arrêt d'urgence &
\textcolor{vespagreen}{matériel réel} \\
Protocole CAN & Sérialisation, priorités, \emph{heartbeats} &
\textcolor{vespaaccent}{\texttt{vcan0} seulement} \\
\midrule
\textbf{Bus CAN physique} & \textcolor{vespared}{\textbf{Ne fonctionne pas}} — \emph{pinmux}
non réclamé & — \\
\textbf{Détection du frelon} & \textcolor{vespared}{\textbf{Absente}} — remplacée par ArUco
& — \\
\textbf{Watchdog} & \textcolor{vespared}{\textbf{Non implémenté}} & — \\
\textbf{Tir laser} & \textcolor{vespared}{\textbf{Jamais effectué}} (décision de sécurité)
& — \\
\bottomrule
\end{tabularx}
\end{center}

\begin{avertissement}[Avant toute mise sous tension du laser]
Le laser est de \textbf{Classe 4}. Il n'a jamais été alimenté, faute de watchdog. Lire
\fichier{docs/annexes/securite\_laser.md} \textbf{en entier} avant d'y toucher. La
collimation dynamique réduit le risque d'incendie, \textbf{pas} le risque oculaire, qui
persiste sur près de 200~m.
\end{avertissement}

% =====================================================================
\newpage
\bibliographystyle{plain}
\bibliography{vespa}

\end{document}
